Bienvenido a la plantilla de \LaTeX para trabajos de graduación de la Universidad del Valle de Guatemala. El objetivo principal de la misma es facilitar la redacción del contenido del trabajo, evitando los posibles problemas que puedan surgir con el formato. En cada sección se incluye una breve descripción sobre qué debe contener la misma (esto corresponde a la información contenida en la guía oficial de la universidad para la redacción de trabajos de graduación), y más adelante se explicará el uso de la plantilla como tal.

\noindent\rule{\linewidth}{0.4pt}

El prefacio es opcional y se coloca antes del índice general. Es la primera que lleva numeración de página con números romanos. Aquí se da las gracias a colaboradores, familia, amigos, etc.; por lo tanto, no debe agregarse un capítulo de agradecimientos. 

En este capítulo, se agradece a las personas o entidades que apoyaron al trabajo; \textbf{no es la introducción}.